\chapter*{Résumé / Abstract}

\begin{otherlanguage}{french}
\section*{Résumé}
Ce rapport présente une étude théorique et numérique de la propagation d'un paquet d'onde au sein d'une chaîne 1D constituée de masses et de ressorts, soumise aux effets combinés du désordre et des non-linéarités. L'analyse quantifie l'impact des fluctuations stochastiques des masses ($m_i$) et des raideurs ($k_i$) sur le transport ondulatoire. Nous caractérisons la transition du régime de propagation balistique dans le régime de localisation forte, où l'énergie subit un confinement spatial exponentiel. Dans un premier temps, on rappelle le formalisme des matrices de transfert, l'exposant de Lyapunov et le rapport de participation pour évaluer la localisation des modes propres. Les résultats numériques sont confrontés aux prédictions analytiques issues de la theorie perturbative.

\vspace*{2em}

\textbf{Mots-clés :} Chaînes harmoniques 1D, Transport ondulatoire, Exposant de Lyapunov, Matrices de transfert.
\end{otherlanguage}

\vspace{1cm}

\begin{otherlanguage}{english}
\section*{Abstract}
This report presents a theoretical and numerical study of wave packet propagation within a 1D chain consisting of masses and springs, subject to the combined effects of disorder and non-linearities. The analysis quantifies the impact of stochastic fluctuations in masses ($m_i$) and stiffnesses ($k_i$) on wave transport. We characterize the transition from the ballistic propagation regime to the strong localization regime, where energy undergoes exponential spatial confinement.First, we recall the transfer matrix formalism, the Lyapunov exponent, and the participation ratio to assess the localization of eigenmodes. Numerical results are compared with analytical predictions from perturbative theory.

\vspace*{2em}

\textbf{Keywords :} 1D harmonic chains, Wave transport, Lyapunov exponent, Transfer matrix.
\end{otherlanguage}